\subsection{Eventual block-intersection consequences}\label{subsec:block_intersection_eventual}

\begin{lemma}[Period-window form of the block intersection identity]
    \label{lem:period_window_block_restriction_intersection}
    \lean{MPSTensor.pgvwc07_iSup_restriction_intersection_eventually_of_period_window}
    \leanok
    \uses{lem:period_window_word_tuple_span_padding,
        lem:block_restriction_intersection_subspace}
    Let \(p>0\). Suppose the length-\(p\) simultaneous block-word tuples span
    \(\prod_j M_{D_j}(\mathbb C)\), and suppose that for every \(0\le r<p\),
    the length-\(s+r\) simultaneous block-word tuples span
    \(\prod_j M_{D_j}(\mathbb C)\). Assume also the normalization
    \[
        \sum_a A^j_aA^{j\,\dagger}_a=I
    \]
    for every block \(j\). Then there is a length \(N\) such that, for every
    \(n\ge N\), if
    \[
        S_n:=\bigvee_jG_{n+1}(A^j),
    \]
    then
    \[
        \left\{\psi:\forall b,\ \psi(-,b)\in S_n\right\}
        \cap
        \left\{\psi:\forall a,\ \psi(a,-)\in S_n\right\}
        =
        \bigvee_jG_{n+2}(A^j).
    \]
\end{lemma}

\begin{proof}
    \leanok
    Lemma~\ref{lem:period_window_word_tuple_span_padding} gives a length \(N\)
    such that the simultaneous block-word tuples span
    \(\prod_j M_{D_j}(\mathbb C)\) at every length \(n\ge N\). For each such
    \(n\), Lemma~\ref{lem:block_restriction_intersection_subspace} applies and
    gives the displayed subspace equality.
\end{proof}

\begin{lemma}[BNT block separation gives one block-intersection step]
    \label{lem:bnt_direct_sum_block_separation_block_intersection}
    \lean{MPSTensor.pgvwc07_iSup_restriction_intersection_of_bnt_directSum_selectors_c1}
    \leanok
    \uses{lem:bnt_direct_sum_block_separation_word_span,
        lem:block_restriction_intersection_subspace}
    Under the hypotheses of
    Lemma~\ref{lem:bnt_direct_sum_block_separation_word_span}, assume in addition the
    normalization
    \[
        \sum_a A^j_aA^{j\,\dagger}_a=I
    \]
    for every block \(j\). Let
    \[
        n=(L_0+1)+(r-1)3(L_0+1),
        \qquad
        S_n=\bigvee_jG_{n+1}(A^j).
    \]
    Then
    \[
        \left\{\psi:\forall b,\ \psi(-,b)\in S_n\right\}
        \cap
        \left\{\psi:\forall a,\ \psi(a,-)\in S_n\right\}
        =
        \bigvee_jG_{n+2}(A^j).
    \]
\end{lemma}

\begin{proof}
    \leanok
    Lemma~\ref{lem:bnt_direct_sum_block_separation_word_span} gives
    \[
        \operatorname{span}\{(A^1_w,\ldots,A^r_w): |w|=n\}
        =
        \prod_jM_{D_j}(\mathbb C).
    \]
    Substituting this span into
    Lemma~\ref{lem:block_restriction_intersection_subspace}, together with the
    normalization, gives the stated subspace equality.
\end{proof}

\begin{lemma}[BNT block separation with a block-injectivity window]
    \label{lem:bnt_direct_sum_block_separation_period_window_block_intersection}
    \lean{MPSTensor.pgvwc07_iSup_restriction_intersection_eventually_of_bnt_directSum_period_window}
    \leanok
    \uses{lem:direct_sum_two_block_dimension_step,
        lem:pair_trace_separation_to_pairwise_block_separation,
        lem:block_separating_equations_from_pairwise_separation,
        lem:period_window_block_restriction_intersection}
    Let \(A^1,\ldots,A^r\) be separated BNT blocks satisfying the
    irreducibility, left-canonicality, and normalized self-overlap hypotheses
    used in the direct-sum comparison. Assume \(L>1\), and assume that, for
    every block \(j\),
    \[
        \operatorname{span}\{A^j_a:a=0,\ldots,d-1\}=M_{D_j}(\mathbb C),
        \qquad
        \operatorname{span}\{A^j_u:|u|=L\}=M_{D_j}(\mathbb C),
        \qquad
        \operatorname{span}\{A^j_v:|v|=L+(L+L)\}=M_{D_j}(\mathbb C).
    \]
    Set
    \[
        S=L+(L+L),
        \qquad
        q=(r-1)S.
    \]
    Suppose that \(p>0\) and let \(s_0\) be a starting length. Assume that, for
    every block \(j\),
    \[
        \operatorname{span}\{A^j_u:|u|=p\}=M_{D_j}(\mathbb C),
    \]
    and that, for every \(0\le s<p+q\) and every block \(j\),
    \[
        \operatorname{span}\{A^j_v:|v|=s_0+s\}=M_{D_j}(\mathbb C).
    \]
    Assume also the normalization
    \[
        \sum_a A^j_aA^{j\,\dagger}_a=I
    \]
    for every block \(j\). Then there is \(N\) such that, for every \(n\ge N\),
    if
    \[
        S_n:=\bigvee_jG_{n+1}(A^j),
    \]
    then
    \[
        \left\{\psi:\forall b,\ \psi(-,b)\in S_n\right\}
        \cap
        \left\{\psi:\forall a,\ \psi(a,-)\in S_n\right\}
        =
        \bigvee_jG_{n+2}(A^j).
    \]
\end{lemma}

\begin{proof}
    \leanok
    The direct-sum dimension step gives, for each ordered pair of distinct
    blocks, length-\(S\) equations which are the identity on the first block and
    zero on the second. Multiplying these equations over the \(r-1\) other
    blocks gives a block-separating equation of total length \(q\). Combining
    this equation with block injectivity at length \(p\) gives the full product
    span at length \(p+q\). Combining it with the assumed block-injectivity
    window gives the full product span at lengths \(s_0+q+s\), for
    \(0\le s<p+q\). The period-window form of the block-intersection identity
    then gives the displayed equality for all sufficiently large \(n\).
\end{proof}

\begin{lemma}[Normalized block-separating equations give all large product spans]
    \label{lem:unital_block_separating_product_span}
    \lean{MPSTensor.wordTupleSpanTop_of_ge_of_common_blockInjective_of_unital_of_pairBlockSeparatingWords}
    \lean{MPSTensor.wordTupleSpanTop_of_ge_of_bnt_directSum_unital_c1}
    \leanok
    \uses{def:pairwise_block_separating_words,
        lem:bnt_direct_sum_block_separation_word_span,
        thm:wordspan_propagation_unital}
    Let \(A^1,\ldots,A^r\) be blocks satisfying the normalization
    \[
        \sum_a A^j_aA^{j\,\dagger}_a=I
    \]
    and assume that every block is injective at length \(L\):
    \[
        \operatorname{span}\{A^j_u:|u|=L\}=M_{D_j}(\mathbb C).
    \]
    If there are block-separating equations of length \(S\), then, for every
    \(n\ge L+(r-1)S\),
    \[
        \operatorname{span}\{(A^1_w,\ldots,A^r_w):|w|=n\}
        =
        \prod_j M_{D_j}(\mathbb C).
    \]
    In particular, for separated BNT blocks satisfying the same normalization,
    if \(L_0>0\) and \(\Gamma_{L_0}^{A^j}\) is injective for every block \(j\),
    the same conclusion holds with
    \[
        S=3(L_0+1),
        \qquad
        n\ge (L_0+1)+(r-1)3(L_0+1).
    \]
\end{lemma}

\begin{proof}
    \leanok
    By Theorem~\ref{thm:wordspan_propagation_unital}, the normalization
    propagates the length-\(L\) span equality to every larger homogeneous length:
    \[
        m\ge L
        \quad\Longrightarrow\quad
        \operatorname{span}\{A^j_u:|u|=m\}=M_{D_j}(\mathbb C).
    \]
    For \(n\ge L+(r-1)S\), take \(m=n-(r-1)S\). Combining the length-\(m\)
    block span with the \(r-1\) block-separating equations gives the displayed
    product span at length \(n\). In the BNT case, the injectivity of
    \(\Gamma_{L_0}^{A^j}\) gives the full span at length \(L_0\), and the
    normalization propagates it to the lengths \(L_0+1\) and \(3(L_0+1)\).
    Lemma~\ref{lem:bnt_direct_sum_block_separation_word_span} then supplies the
    block-separating equations with \(S=3(L_0+1)\).
\end{proof}

\begin{lemma}[Product span makes the sum of local spaces direct]
    \label{lem:product_span_block_image_direct_sum}
    \lean{MPSTensor.groundSpace_iSupIndep_of_wordTupleSpanTop}
    \lean{MPSTensor.groundSpace_iSupIndep_of_ge_of_bnt_directSum_unital_c1}
    \leanok
    \uses{def:ground_space_parent, lem:unital_block_separating_product_span}
    Suppose
    \[
        \operatorname{span}\{(A^1_w,\ldots,A^r_w):|w|=n\}
        =
        \prod_jM_{D_j}(\mathbb C).
    \]
    Then the sum of the local spaces \(G_n(A^j)\) is direct: if
    \[
        \phi_j\in G_n(A^j),
        \qquad
        \sum_j\phi_j=0,
    \]
    then \(\phi_j=0\) for every \(j\).  In particular, for the normalized BNT
    hypotheses above, this directness holds for every
    \[
        n\ge (L_0+1)+(r-1)3(L_0+1).
    \]
\end{lemma}

\begin{proof}
    \leanok
    Write
    \[
        \phi_j(\sigma)=\tr(A^j_\sigma X_j).
    \]
    Evaluating \(\sum_j\phi_j=0\) on each word \(w\) of length \(n\) gives
    \[
        \sum_j\tr(A^j_wX_j)=0.
    \]
    Since \(\tr(A^j_wX_j)=\tr(X_jA^j_w)\), the product-span equation gives, for
    every \((M_1,\ldots,M_r)\in\prod_jM_{D_j}(\mathbb C)\),
    \[
        \sum_j\tr(X_jM_j)=0.
    \]
    Taking only the \(j\)-th component nonzero and using nondegeneracy of the
    trace pairing yields \(X_j=0\), hence \(\phi_j=0\). The BNT statement follows
    by substituting Lemma~\ref{lem:unital_block_separating_product_span}.
\end{proof}
